\documentclass[11pt]{article}
\usepackage[margin=1.2in, top=1.0in, bottom=1.0in]{geometry}
\usepackage{amsmath, amssymb, amsfonts}
\usepackage{microtype}
\usepackage{xspace}
\usepackage[T1]{fontenc}
\usepackage{lmodern}
\usepackage{parskip}

% --- Macros ---
\newcommand{\method}{\textsc{RetroDiff}\xspace}
\newcommand{\Lobj}{\mathcal{L}}
\newcommand{\Lelbo}{\mathcal{L}_{\mathrm{ELBO}}}
\newcommand{\Lretro}{\mathcal{L}_{\mathrm{RetroDiff}}}
\newcommand{\KL}{\mathrm{KL}}
\newcommand{\Ret}{R(x)}
\newcommand{\ptheta}{p_\theta}
\newcommand{\Eq}[1]{\mathbb{E}_{#1}}

\pagestyle{empty}

\begin{document}

\begin{center}
  {\large\bfseries Training Objective Derivation for \method}
\end{center}

\vspace{0.5em}

\noindent
The \method training objective follows the standard evidence lower bound (ELBO) framework
for latent-variable diffusion models, then extends it to the retrieval-augmented setting.
We derive the three-step chain below.

\medskip
\noindent\textbf{Step 1: Variational lower bound.}
Let $q(x_{1:T}\mid x)$ be the fixed forward (noising) process that gradually corrupts a
clean sample $x$ over $T$ steps, and let $\ptheta(x_{0:T})$ be the learned joint
generative model.  Because the log-likelihood $\log \ptheta(x)$ is intractable, we bound
it from below via Jensen's inequality applied to the concave $\log$:

\begin{align}
  \log \ptheta(x)
    \;\geq\;
    \Eq{q(x_{1:T}|x)}\!\left[\log\frac{\ptheta(x_{0:T})}{q(x_{1:T}|x)}\right]
    \;=\;
    -\Lelbo
  \label{eq:elbo}
\end{align}

\noindent
The inequality is tight when $q$ matches the true posterior $\ptheta(x_{1:T}\mid x_0)$;
in practice optimising $-\Lelbo$ drives the model towards high likelihood without
evaluating an intractable normaliser.

\medskip
\noindent\textbf{Step 2: Reconstruction--KL decomposition.}
Expanding the ELBO by Markov factorisation of both $\ptheta$ and $q$ along the diffusion
chain yields a reconstruction term at $t\!=\!1$ plus one KL divergence per denoising step:

\begin{align}
  \Lelbo
  \;=\;
  \underbrace{\Eq{q}\bigl[-\log \ptheta(x_0\mid x_1)\bigr]}_{\text{reconstruction}}
  \;+\;
  \sum_{t=2}^{T}
    \underbrace{\KL\!\left(q(x_{t-1}\mid x_t, x)
               \,\big\|\,
               \ptheta(x_{t-1}\mid x_t)\right)}_{\text{step-}t\text{ denoising}}
  \label{eq:decomp}
\end{align}

\noindent
Each KL term measures how well the learned reverse step $\ptheta(x_{t-1}\mid x_t)$ matches
the tractable forward posterior $q(x_{t-1}\mid x_t, x)$.  Because the forward process has
Gaussian transitions, both quantities are Gaussian and each KL reduces to a weighted
$\ell_2$ loss on the predicted noise, recoverable in closed form.

\medskip
\noindent\textbf{Step 3: Retrieval-augmented objective (\method contribution).}
Standard diffusion trains the reverse process without side information.
\method augments each sample $x$ with a set $\Ret$ of $k$ retrieved exemplars drawn from a
nearest-neighbour index over the training corpus.  The per-sample ELBO is conditioned on
these exemplars, and the outer expectation averages over the data distribution
$p_{\mathrm{data}}$:

\begin{align}
  \Lretro(\theta)
  \;=\;
  \Eq{x \sim p_{\mathrm{data}}}\!\left[
      \Lelbo\!\left(x \;\big|\; \Ret\right)
  \right]
  \label{eq:retro}
\end{align}

\noindent
Conditioning the ELBO on $\Ret$ propagates retrieved appearance into every denoising KL
term in~\eqref{eq:decomp}: each reverse step $\ptheta(x_{t-1}\mid x_t, \Ret)$ now has
access to exemplar context.  The reconstruction term in~\eqref{eq:elbo} similarly becomes
$-\log\ptheta(x_0\mid x_1, \Ret)$.  Minimising~\eqref{eq:retro} therefore encourages the
model to exploit retrieved structure across the entire denoising chain, rather than at a
single conditioning step.

\end{document}
